\section{前期调研与方法选择}
\label{sec:related}
\subsection{传统匿名化方法}
调研首先比较模糊、马赛克和遮挡。它们不需要额外训练，计算量小，适合建立基础对照。比较时拟使用同一批带标签的人脸，观察编辑前后眼眉口结构的可见程度，并统计匿名图表情准确率和身份预测变化。这样可以判断简单处理是否足以支撑后续识别，而不是只凭视觉上“看不清脸”选型。

这类方法的处理区域与表情关键区域重叠，强度增加时两类信息会一起减弱。因此，拟将其作为基础匿名方式进行分析，并进一步调研能够分区域控制的路线。

\subsection{整脸替换与生成式方法}
整脸替换和生成式方法通过改变外观降低身份关联，比较重点拟放在三个方面：表情与姿态是否跟随源图，合成是否改动头发和背景，相邻帧之间是否连贯。项目将对同一输入分别观察五官形态、边界过渡和处理耗时，再判断其与车载演示要求的匹配程度。

本项目要求保留原场景，并只对面部内部进行编辑。因此，完整替换整张人脸不作为首选，生成式组件更适合在明确的局部掩膜和几何条件下使用。

\subsection{局部编辑与关键点约束}
局部残差能够在输入基础上控制变化幅度，关键点能够提供眼睛开合、眉毛位置与嘴部形状的约束。GANonymization 的关键点条件生成思路为这一选择提供参考\cite{ganonymization}。拟把局部编辑、几何约束与关键点条件组织为统一模块，再通过掩膜合成回灰度画面。

\begin{table}[htbp]
\centering
\caption{备选方案的比较重点与选型考虑}
\label{tab:method_compare}
\begin{tabularx}{\linewidth}{@{}lXX@{}}
\toprule
备选路线 & 拟比较的主要内容 & 选型考虑\\
\midrule
模糊、马赛克 & 身份预测变化及表情结构可见性 & 用作基础处理对照\\
整脸替换 & 表情迁移、背景边界与连续性 & 与局部编辑要求比较\\
关键点条件 & 眼眉口几何与局部生成效果 & 作为表情保持条件\\
局部残差 & 修改范围、变化强度及处理开销 & 拟作为主体处理方式\\
\bottomrule
\end{tabularx}
\end{table}

综合上述比较，项目拟采用局部残差与关键点约束相结合的方案。选定的理由是修改范围可控、能够保留动作关系，并便于与轻量表情模型衔接。实际参数将在训练集上调整，候选方案在固定验证集上比较后确定，不采用先预测原图表情再选择替换图片的流程。
